% !PRINTED 21
\weberpage{21}
%<<CONT
be. These theorems hold if only one of the two groups $Q$, $Q'$
is a greatest normal divisor of $P$. We now assume that they
both have this property, and understand by $R$ the greatest
common divisor of $Q$ and $Q'$. This group $R$ is then a normal
divisor of $P$ as well as of $Q$ and $Q'$. For if $a$ is any
element in $P$, and $c$ is contained in $Q$ as well as in $Q'$,
then $a^{-1}ca$ is also in $Q$ and in $Q'$, and hence also
contained in $R$. Therefore $a^{-1}Ra = R$, and $R$ is a normal
divisor of $P$ and consequently a normal divisor of $Q$ and of $Q'$
(\S.\,4, 5.). In order to decompose $Q$ into the subsidiary groups
of $R$, we choose a suitable system of elements $1, b_1, b_2, b_3 \ldots$
in $Q$, so that $R, Rb_1, Rb_2, Rb_3 \ldots$ all become different
from one another, and set
\[
Q = R + Rb_1 + Rb_2 + Rb_3 + \ldots,
\]
which we can also, if we set
\weq{18}{B = 1, b_1, b_2, b_3 \ldots}
denote briefly by
\weq{19}{Q = RB}
Now by (13) we have $Q'Q = P$, hence also
\weq{20}{P = Q'RB = Q'B}
or
\[
P = Q' + Q'b_1 + Q'b_2 + Q'b_3 \ldots
\]

The subsidiary groups
\weq{21}{Q',\ Q'b_1,\ Q'b_2,\ Q'b_3 \ldots}
are, however, all different from one another. For if we had, say,
\[
Q'b_2 = Q'b_1,
\]
it would follow, since $Q'$ contains the unit, that there is an element
$e$ in $Q'$ such that
\[
b_2 = eb_1
\]
would hold. This element $e$, however, is equal to $b_2b_1^{-1}$, and
hence, since the $b$ are contained in $Q$, is likewise contained in $Q$
and therefore also in $R$. But then we should also have $Rb_2 = Rb_1$,
which is contrary to the assumption.

Accordingly we can form the elements of the group complementary to $R$
with respect to $Q$ and of the group complementary to $Q'$ with respect
to $P$. We obtain these two groups
\weq{22}{\begin{aligned}
Q/R &= R,\ Rb_1,\ Rb_2,\ Rb_3 \ldots \\
P/Q' &= Q',\ Q'b_1,\ Q'b_2,\ Q'b_3 \ldots
\end{aligned}}
